% Knowledge dependencies (solid) and feedback iteration (dashed).
\begin{tikzpicture}[
  x=1mm,y=1mm,font=\figurefont\fontsize{8}{10}\selectfont,
  box/.style={draw=BookRule,fill=BookPaper,minimum height=11mm,
    text width=27mm,align=center,inner sep=2mm,line width=.7pt},
  flow/.style={-{Stealth[length=2mm]},draw=BookMuted,line width=.8pt},
  loop/.style={flow,dashed},
  lab/.style={fill=white,inner sep=1mm,font=\figurefont\fontsize{7}{9}\selectfont}
]
  \node[box,fill=FigureInput!10] (limits) at (16,28) {三类奖励限制};
  \node[box] (obs) at (64,51) {观测\\可推断什么};
  \node[box] (credit) at (64,28) {归因\\如何分配结果};
  \node[box] (goal) at (64,5) {目标匹配\\行为是否符合要求};
  \node[box,fill=FigureModel!10] (update) at (112,39) {整合与更新};
  \node[box,fill=FigureOutput!10] (check) at (112,5) {独立检验\\继续或停止};
  \node[box,text width=20mm] (query) at (153,39) {定向查询\\新反馈};
  \draw[flow] (limits.east) -- ++(5,0) |- (obs.west);
  \draw[flow] (limits) -- (credit);
  \draw[flow] (limits.east) -- ++(5,0) |- (goal.west);
  \draw[flow] (obs.east) -- ++(6,0) |- (update.west);
  \draw[flow] (credit.east) -- ++(6,0) |- (update.west);
  \draw[flow] (goal) -- (check);
  \draw[loop] (update) -- node[lab,right] {新行为} (check);
  \draw[loop] (check.east) -- (153,5)
    -- node[lab,right,align=left] {有可补缺口\\且有预算} (query.south);
  \draw[loop] (query) -- node[lab,above] {反馈} (update);
  \draw[loop] (check.south) -- ++(0,-10)
    node[below,align=center] {满足要求／预算耗尽／识别受限\\分别记录停止结论};
\end{tikzpicture}
